\documentclass[11pt,a4paper]{article}

\usepackage[T1]{fontenc}
\usepackage[utf8]{inputenc}
\usepackage{lmodern}
\usepackage{amsmath,amssymb,amsthm,mathtools}
\usepackage{longtable}
\usepackage{booktabs}
\usepackage{enumitem}
\usepackage{fancyhdr}
\usepackage[margin=1in,bottom=1.5in,footskip=0.5in,headheight=14pt]{geometry}
\usepackage[colorlinks=true,linkcolor=blue,citecolor=blue,urlcolor=blue]{hyperref}
\usepackage[nameinlink,capitalise]{cleveref}

\theoremstyle{plain}
\newtheorem{fact}{Fact}
\newtheorem{lemma}{Lemma}
\newtheorem{proposition}{Proposition}
\newtheorem{theorem}{Theorem}
\newtheorem{corollary}{Corollary}
\theoremstyle{definition}
\newtheorem{assumption}{Assumption}
\newtheorem{remark}{Remark}

\crefname{fact}{Fact}{Facts}

\newcommand{\R}{\mathbb{R}}
\newcommand{\V}{\mathcal{V}}
\newcommand{\E}{\mathcal{E}}
\newcommand{\G}{\mathcal{G}}
\DeclareMathOperator{\sgn}{sgn}
\DeclareMathOperator{\dist}{dist}
\DeclareMathOperator*{\argmax}{arg\,max}
\DeclarePairedDelimiter{\norm}{\lVert}{\rVert}
\DeclarePairedDelimiter{\abs}{\lvert}{\rvert}

% ---------------------------------------------------------------------------
% Per-page notation key. Each section declares the active context; the footer
% prints only the symbols that appear on the current page range.
% ---------------------------------------------------------------------------
\newif\ifnotationkey
\notationkeytrue

\makeatletter
\newcommand{\DeclareNotationKey}[2]{%
  \expandafter\gdef\csname nmkey@#1\endcsname{#2}%
}
\newcommand{\NotationContext}[1]{\markboth{#1}{#1}}
\newcommand{\nmkeyCurrent}{%
  \edef\nmkey@id{\firstmark}%
  \ifx\nmkey@id\@empty
    \edef\nmkey@id{\rightmark}%
  \fi
  \ifx\nmkey@id\@empty\else
    \@ifundefined{nmkey@\nmkey@id}{}{%
      \csname nmkey@\nmkey@id\endcsname
    }%
  \fi
}
\newcommand{\nmkeyFooter}{%
  \ifnotationkey
    \footnotesize
    \begin{minipage}{\textwidth}
      \rule{\textwidth}{0.4pt}\\[2pt]
      \textsc{key}\quad\nmkeyCurrent
    \end{minipage}%
  \fi
}
\makeatother

\pagestyle{fancy}
\fancyhf{}
\renewcommand{\headrulewidth}{0pt}
\renewcommand{\footrulewidth}{0pt}
\fancyfoot[C]{\nmkeyFooter}
\renewcommand{\sectionmark}[1]{}
\renewcommand{\subsectionmark}[1]{}

\DeclareNotationKey{facts}{%
$n$ robots, ring radius $R$, centroid $p^*$, source $p_s$, error $e=p^*-p_s$.
$\phi_i$ unit slot direction, $z_i=p_i-R\phi_i$ shifted state, $q_i=z_i-\bar z$ disagreement.
$\alpha$ sign gain, $\beta$ measurement gain, $\kappa$ field curvature, $\delta$ noise bound,
$\lambda=2\beta\kappa$, $\varepsilon$ ultimate error radius.}

\DeclareNotationKey{teams}{%
$N$ sources, $p_s^{(k)}$ source $k$, $\mathcal{B}_k$ its Voronoi basin, $d_k$ source-to-seam clearance.
Team $k$: index set $\V_k$, size $n_k$, radius $R_k$, centroid $p_k^*$, informed fraction $\chi_k$,
formation deviation $q_k$, error radius $\varepsilon_k$, split time $t_s$.}

\DeclareNotationKey{moving}{%
$p_s(t)$ moving source, $v_{\max}$ speed bound, $v$ constant velocity,
$e=p^*-p_s$, $\lambda=2\beta\kappa$, $d_\eta$ noise-induced drift, $T_f$ formation time,
$\hat v$ velocity feedforward, $\tilde v=\hat v-\dot p_s$.}

\DeclareNotationKey{field}{%
$\Omega$ search domain, $f$ nonconvex field, $\ell_j$ local maximisers ordered by value,
$f^\star=f(\ell_1)$ global value, $\Delta$ value gap, $L_f$ Lipschitz constant of $f$,
$L_H$ Lipschitz constant of $\nabla^2f$, $g_R$ circular gradient estimate, $r_R$ its bias.}

\DeclareNotationKey{multistart}{%
$M$ teams, $\hat\ell_k$ terminal centroid of team $k$, $\hat F_k$ its reported score,
$\varepsilon_{\mathrm{loc}}$ terminal localisation error, $\delta_F$ score error,
$\Delta$ value gap, $L_f$ Lipschitz constant of $f$.}

\allowdisplaybreaks

\begin{document}

\section{Standing results used without proof}\label{sec:facts}
\NotationContext{facts}

The three gap sections below are self-contained except for five properties of the
single-source baseline of \cite{du2024}. They are collected here as facts and are not
re-derived; every other statement in this document is proved step by step.

Throughout, $n$ robots with positions $p_i\in\R^2$ communicate over an undirected graph
$\G=(\V,\E)$, are assigned distinct slot directions
$\phi_i=(\cos\theta_i,\sin\theta_i)$ with $\theta_i=2\pi(i-1)/n$, and run
\begin{equation}\label{eq:baseline-controller}
    u_i=\alpha\sum_{j\in\mathcal{N}_i}\sgn\!\left(z_j-z_i\right)
       -\frac{2\beta}{R}\,\sigma_i\phi_i,
    \qquad z_i\coloneqq p_i-R\phi_i,
\end{equation}
where $\sigma_i$ is robot $i$'s scalar measurement, $\sgn$ acts componentwise, and
$\mathcal{N}_i$ is the neighbour set. The centroid is $p^*\coloneqq n^{-1}\sum_ip_i$.

\begin{fact}[Slot-direction moment identities]\label{fact:moments}
For $\theta_i=2\pi(i-1)/n$ and $n\geq3$,
\begin{equation}\label{eq:moments}
    \sum_{i=1}^n\phi_i=0,
    \qquad
    \sum_{i=1}^n\phi_i\phi_i^{\top}=\frac n2 I_2,
\end{equation}
and if in addition $n\geq4$ then for every symmetric $H\in\R^{2\times2}$
\begin{equation}\label{eq:third-moment}
    \sum_{i=1}^n\left(\phi_i^{\top}H\phi_i\right)\phi_i=0.
\end{equation}
The restriction $n\geq4$ in \eqref{eq:third-moment} is essential and is the source of the
$n=3$ anomaly analysed in \cref{sec:gap3-bias}.
\end{fact}

\begin{fact}[Sign-dissipation and disagreement comparison]\label{fact:sign}
Let $q_i\coloneqq z_i-n^{-1}\sum_jz_j$, $V_f\coloneqq\frac12\sum_i\norm{q_i}^2$,
$S(z)\coloneqq\sum_{(i,j)\in\E}\norm{z_i-z_j}_1$ and $Q(z)\coloneqq\sum_i\norm{q_i}$.
If $\G$ is connected then
\begin{equation}\label{eq:sign-dissipation}
    \sum_i q_i^{\top}\sum_{j\in\mathcal{N}_i}\sgn\!\left(z_j-z_i\right)=-S(z),
    \qquad
    Q(z)\leq n\,S(z),
    \qquad
    S(z)\geq\frac{\sqrt{2V_f}}{n}.
\end{equation}
\end{fact}

\begin{fact}[Finite-time formation]\label{fact:formation}
Assume $\G$ connected, $n\geq3$, and $\abs{\sigma_i}\leq f_{D_{\max}}$ for all $i$ and all $t$.
If the gain condition
\begin{equation}\label{eq:gain-condition}
    \gamma\coloneqq\alpha-\frac{2\beta n f_{D_{\max}}}{R}>0
\end{equation}
holds, then $V_f(t)=0$ for all $t\geq T_f$ with
\begin{equation}\label{eq:formation-time}
    T_f\leq\frac{\sqrt2\,n\sqrt{V_f(0)}}{\gamma}.
\end{equation}
\end{fact}

\begin{fact}[All-informed localisation]\label{fact:localisation}
Assume the quadratic bowl $f(p)=\kappa\norm{p-p_s}^2$ with $\kappa>0$, every robot informed
with $\sigma_i=f(p_i)+\eta_i$ and $\abs{\eta_i}\leq\delta$, and exact formation. Then with
$e\coloneqq p^*-p_s$ and $\lambda\coloneqq2\beta\kappa$,
\begin{equation}\label{eq:baseline-error-dynamics}
    \dot e=-\lambda e+d_\eta,
    \qquad
    d_\eta=-\frac{2\beta}{nR}\sum_i\eta_i\phi_i,
    \qquad
    \norm{d_\eta}\leq\frac{2\beta\delta}{R},
\end{equation}
and consequently $\limsup_{t\to\infty}\norm{e(t)}\leq\delta/(\kappa R)$.
\end{fact}

\begin{fact}[Partial-information error radius]\label{fact:epsilon}
Under the saturation geometry of \cite{du2024}, if at least $n_{\mathcal X}$ of the $n$ robots
are informed at every time, the ultimate centroid error obeys
$\limsup_{t\to\infty}\norm{e(t)}\leq\varepsilon(n_{\mathcal X})$ with
\begin{equation}\label{eq:du-epsilon}
    \varepsilon(n_{\mathcal X})
    =\frac{2\pi n\delta}
    {\kappa R\left(2\pi n_{\mathcal X}-n\abs{\sin\!\left(2\pi n_{\mathcal X}/n\right)}\right)}.
\end{equation}
\end{fact}

\begin{lemma}[The error radius is always well posed]\label{lem:denominator}
For every $n\geq3$ and every $1\leq n_{\mathcal X}\leq n$ the denominator in
\eqref{eq:du-epsilon} is strictly positive, so $\varepsilon(n_{\mathcal X})$ is finite and positive.
\end{lemma}

\begin{proof}
Write $x\coloneqq2\pi n_{\mathcal X}/n$, so that $n_{\mathcal X}=nx/(2\pi)$ and $x\in(0,2\pi]$.
The denominator factor becomes
\begin{equation*}
    2\pi n_{\mathcal X}-n\abs{\sin x}
    =2\pi\cdot\frac{nx}{2\pi}-n\abs{\sin x}
    =n\left(x-\abs{\sin x}\right).
\end{equation*}
For $x\geq0$ the mean value theorem applied to $s\mapsto\sin s$ on $[0,x]$ gives
$\abs{\sin x}=\abs{\sin x-\sin0}\leq x\sup_{s}\abs{\cos s}\leq x$, with equality only at $x=0$
because $\abs{\cos s}<1$ on any interval of positive length not reducing to $\{0\}$ modulo the
isolated points where $\abs{\cos s}=1$. Since $x>0$ here, $x-\abs{\sin x}>0$ and the factor is
$n(x-\abs{\sin x})>0$.
\end{proof}

\begin{remark}\label{rem:denominator-caveat}
\cref{lem:denominator} removes a caveat that is usually attached to \eqref{eq:du-epsilon}.
The denominator looks as though it could change sign for small informed counts, and it is often
guarded by an explicit positivity assumption. The change of variable above shows the guard is
vacuous: positivity is automatic for every admissible pair $(n,n_{\mathcal X})$. A direct sweep
over $3\leq n\leq60$ and $1\leq n_{\mathcal X}\leq n$ finds the minimum of the denominator factor
to be $1.15\times10^{-2}$, attained at $(n,n_{\mathcal X})=(60,1)$, consistent with the bound.
What does degrade for small $n_{\mathcal X}$ is the \emph{size} of $\varepsilon$, not its
existence.
\end{remark}

\section{Gap 1: $N$ sources served by $N$ teams}\label{sec:gap1}
\NotationContext{teams}

\subsection{Setting}\label{sec:gap1-setting}

Let $p_s^{(1)},\ldots,p_s^{(N)}\in\R^2$ be distinct sources of equal curvature $\kappa>0$ and let
the sensed field be the lower envelope of the individual bowls,
\begin{equation}\label{eq:multisource-field}
    f(p)=\min_{1\leq k\leq N}f_k(p),
    \qquad f_k(p)=\kappa\norm{p-p_s^{(k)}}^2.
\end{equation}
The basins of \eqref{eq:multisource-field} are the Voronoi cells of the source set
\cite{cortes2004},
\begin{equation}\label{eq:voronoi}
    \mathcal{B}_k=\left\{p\in\R^2:\norm{p-p_s^{(k)}}\leq\norm{p-p_s^{(j)}}\ \forall j\neq k\right\}.
\end{equation}
The $n$ robots are partitioned after a split time $t_s$ into teams $\V_1,\ldots,\V_N$ with
$n_k=\abs{\V_k}$, $\sum_kn_k=n$. Team $k$ has its own ring radius $R_k$, gains
$(\alpha_k,\beta_k)$, slot directions $\phi_i^{(k)}$ indexed within the team, and centroid
$p_k^*=n_k^{-1}\sum_{i\in\V_k}p_i$. Each robot runs \eqref{eq:baseline-controller} with the sum
restricted to $\mathcal{N}_i\cap\V_k$.

Define the source-to-seam clearance and the team formation deviation
\begin{align}
    d_k&\coloneqq\dist\!\left(p_s^{(k)},\partial\mathcal{B}_k\right),\label{eq:clearance-def}\\
    q_k(t)&\coloneqq\max_{i\in\V_k}\norm{p_i(t)-p_k^*(t)-R_k\phi_i^{(k)}}.\label{eq:team-deviation}
\end{align}

\subsection{Step-by-step: the clearance has a closed form}\label{sec:gap1-clearance}

The certificate in \cref{sec:gap1-basin} is stated in terms of $d_k$, which is defined as a
distance to a polygonal boundary. The following lemma reduces it to pairwise source distances.
It is used silently in most treatments; the proof is short but the second half is not immediate,
because a point on the nearest bisector need not lie on $\partial\mathcal{B}_k$ in general.

\begin{lemma}[Clearance equals half the nearest-neighbour distance]\label{lem:clearance}
For the equal-curvature field \eqref{eq:multisource-field},
\begin{equation}\label{eq:clearance-closed-form}
    d_k=\frac12\min_{j\neq k}\norm{p_s^{(k)}-p_s^{(j)}}.
\end{equation}
\end{lemma}

\begin{proof}
Write $\rho\coloneqq\frac12\min_{j\neq k}\norm{p_s^{(k)}-p_s^{(j)}}$ and let $j^\star$ attain the
minimum.

\emph{Step 1 ($d_k\geq\rho$).} Let $y\in\partial\mathcal{B}_k$. By \eqref{eq:voronoi} the
boundary consists of points for which at least one competing constraint is active, so there is
$j\neq k$ with $\norm{y-p_s^{(k)}}=\norm{y-p_s^{(j)}}$. The triangle inequality then gives
\begin{equation*}
    \norm{p_s^{(k)}-p_s^{(j)}}
    \leq\norm{p_s^{(k)}-y}+\norm{y-p_s^{(j)}}
    =2\norm{y-p_s^{(k)}},
\end{equation*}
hence $\norm{y-p_s^{(k)}}\geq\frac12\norm{p_s^{(k)}-p_s^{(j)}}\geq\rho$. Taking the infimum over
$y\in\partial\mathcal{B}_k$ yields $d_k\geq\rho$.

\emph{Step 2 (a boundary point at distance exactly $\rho$).} Let
$m\coloneqq\frac12\left(p_s^{(k)}+p_s^{(j^\star)}\right)$ be the midpoint, so
$\norm{m-p_s^{(k)}}=\norm{m-p_s^{(j^\star)}}=\rho$. For any $l\neq k$,
\begin{equation*}
    \norm{m-p_s^{(l)}}
    \geq\norm{p_s^{(k)}-p_s^{(l)}}-\norm{m-p_s^{(k)}}
    \geq2\rho-\rho=\rho=\norm{m-p_s^{(k)}},
\end{equation*}
where the second inequality uses the definition of $\rho$ as a minimum over all $j\neq k$.
Therefore $m$ satisfies every constraint in \eqref{eq:voronoi} and so $m\in\mathcal{B}_k$; since
the constraint for $j^\star$ holds with equality, $m\in\partial\mathcal{B}_k$. Hence
$d_k\leq\norm{m-p_s^{(k)}}=\rho$.

Combining the two steps gives $d_k=\rho$.
\end{proof}

\subsection{Step-by-step: team size cancels from the error radius}\label{sec:gap1-scaling}

The design question for Gap 1 is how to split $n$ robots into $N$ teams. The following
computation shows that, at the level of the asymptotic error radius, the head count $n_k$ is
\emph{not} the relevant variable.

\begin{lemma}[Team-size cancellation]\label{lem:team-scaling}
Suppose team $k$ satisfies the hypotheses of \cref{fact:epsilon} with $n\to n_k$,
$R\to R_k$, and a guaranteed informed count $n_{\mathcal X,k}=\chi_kn_k$ for some
$\chi_k\in(0,1]$. Then
\begin{equation}\label{eq:team-epsilon-scale}
    \varepsilon_k
    =\frac{2\pi\delta}{\kappa R_k\left(2\pi\chi_k-\abs{\sin\!\left(2\pi\chi_k\right)}\right)},
\end{equation}
which depends on $n_k$ only through the informed \emph{fraction} $\chi_k$. In particular
$\varepsilon_k=\delta/(\kappa R_k)$ when $\chi_k=1$.
\end{lemma}

\begin{proof}
\emph{Step 1 (substitute).} Replacing $n\to n_k$, $R\to R_k$ and
$n_{\mathcal X}\to\chi_kn_k$ in \eqref{eq:du-epsilon},
\begin{equation*}
    \varepsilon_k
    =\frac{2\pi n_k\delta}
    {\kappa R_k\left(2\pi\chi_kn_k-n_k\abs{\sin\!\left(2\pi\chi_kn_k/n_k\right)}\right)}.
\end{equation*}

\emph{Step 2 (simplify the sine argument).} The argument of the sine is
$2\pi\chi_kn_k/n_k=2\pi\chi_k$, which no longer involves $n_k$:
\begin{equation*}
    \varepsilon_k
    =\frac{2\pi n_k\delta}
    {\kappa R_k\left(2\pi\chi_kn_k-n_k\abs{\sin\!\left(2\pi\chi_k\right)}\right)}.
\end{equation*}

\emph{Step 3 (factor).} Both terms in the bracket carry a factor $n_k$:
\begin{equation*}
    \varepsilon_k
    =\frac{2\pi n_k\delta}
    {\kappa R_k\,n_k\left(2\pi\chi_k-\abs{\sin\!\left(2\pi\chi_k\right)}\right)}.
\end{equation*}

\emph{Step 4 (cancel).} By \cref{lem:denominator} with $x=2\pi\chi_k>0$ the bracket is strictly
positive, so cancelling the common factor $n_k\neq0$ is legitimate and yields
\eqref{eq:team-epsilon-scale}.

\emph{Step 5 (all-informed case).} Setting $\chi_k=1$ gives
$\abs{\sin2\pi}=0$, hence
$\varepsilon_k=2\pi\delta/(\kappa R_k\cdot2\pi)=\delta/(\kappa R_k)$, which agrees with
\cref{fact:localisation} as it must.
\end{proof}

\begin{remark}[Consequence for team sizing]\label{rem:team-sizing}
\eqref{eq:team-epsilon-scale} says that adding robots to a team improves its asymptotic accuracy
only insofar as it raises the informed fraction or permits a larger ring radius $R_k$. Two
allocations with the same $(\chi_k,R_k)$ but different $n_k$ have identical predicted error
radii. Integrality is the only caveat: $\chi_kn_k$ must be an integer count, so for small $n_k$
the achievable fractions form a coarse grid, and $\chi_k$ should be read as the guaranteed lower
fraction rather than an instantaneous one.
\end{remark}

\subsection{Step-by-step: basin containment}\label{sec:gap1-basin}

The reduction of the $N$-source problem to $N$ single-source problems is only valid while each
team stays inside its own basin. Centroid membership is not enough: a ring of radius $R_k$ can
straddle a seam even when $p_k^*\in\mathcal{B}_k$. The following certificate is stated on the
robots, not the centroid.

\begin{lemma}[Direct basin-containment certificate]\label{lem:basin-certificate}
Assume \eqref{eq:multisource-field}. If for all $t\geq t_s$
\begin{equation}\label{eq:basin-certificate}
    \norm{p_k^*(t)-p_s^{(k)}}+R_k+q_k(t)<d_k,
\end{equation}
then $p_i(t)\in\operatorname{int}\mathcal{B}_k$ for every $i\in\V_k$ and every $t\geq t_s$.
\end{lemma}

\begin{proof}
Fix $t\geq t_s$ and $i\in\V_k$.

\emph{Step 1 (decompose the offset).} Insert and cancel the nominal slot position:
\begin{equation*}
    p_i-p_s^{(k)}
    =\underbrace{\left(p_i-p_k^*-R_k\phi_i^{(k)}\right)}_{\text{formation error}}
    +\underbrace{R_k\phi_i^{(k)}}_{\text{slot offset}}
    +\underbrace{\left(p_k^*-p_s^{(k)}\right)}_{\text{centroid error}}.
\end{equation*}

\emph{Step 2 (triangle inequality).} Applying it twice,
\begin{equation*}
    \norm{p_i-p_s^{(k)}}
    \leq\norm{p_i-p_k^*-R_k\phi_i^{(k)}}
    +R_k\norm{\phi_i^{(k)}}
    +\norm{p_k^*-p_s^{(k)}}.
\end{equation*}

\emph{Step 3 (evaluate the two known terms).} The slot directions are unit vectors, so
$\norm{\phi_i^{(k)}}=1$. By \eqref{eq:team-deviation} the first term is at most $q_k(t)$, since
$q_k$ is the maximum over $i\in\V_k$ of exactly that quantity. Hence
\begin{equation}\label{eq:basin-step3}
    \norm{p_i-p_s^{(k)}}\leq q_k(t)+R_k+\norm{p_k^*(t)-p_s^{(k)}}.
\end{equation}

\emph{Step 4 (apply the hypothesis).} The right-hand side of \eqref{eq:basin-step3} is the
left-hand side of \eqref{eq:basin-certificate}, so
\begin{equation}\label{eq:basin-step4}
    \norm{p_i-p_s^{(k)}}<d_k.
\end{equation}

\emph{Step 5 (translate into a Voronoi statement).} Let $j\neq k$. By \cref{lem:clearance},
$\norm{p_s^{(k)}-p_s^{(j)}}\geq2d_k$. The reverse triangle inequality and
\eqref{eq:basin-step4} give
\begin{equation*}
    \norm{p_i-p_s^{(j)}}
    \geq\norm{p_s^{(k)}-p_s^{(j)}}-\norm{p_i-p_s^{(k)}}
    >2d_k-d_k
    =d_k
    >\norm{p_i-p_s^{(k)}},
\end{equation*}
where the last inequality is \eqref{eq:basin-step4} again.

\emph{Step 6 (conclude).} The chain holds for every $j\neq k$, so $p_i$ is \emph{strictly}
closer to $p_s^{(k)}$ than to any other source. Every constraint in \eqref{eq:voronoi} is
therefore strict at $p_i$, which is the definition of
$p_i\in\operatorname{int}\mathcal{B}_k$.
\end{proof}

\begin{remark}[Sufficient, not necessary]\label{rem:basin-sufficiency}
\eqref{eq:basin-certificate} charges the full ring radius against the clearance, so it is
conservative: a team can be entirely inside its basin while the certificate fails, for instance
when the ring is elongated away from the seam. What \cref{lem:basin-certificate} rules out is the
converse direction only. This matters for interpretation, because a run in which
\eqref{eq:basin-certificate} fails is not thereby proved unsafe; it is merely uncertified. The
numerical study in \cref{sec:gap1-numerics} exhibits both a certified run and an uncertified run
in which containment genuinely fails, which shows the condition is neither vacuous nor
necessary.
\end{remark}

\subsection{Step-by-step: reduction to $N$ baseline problems}\label{sec:gap1-reduction}

\begin{proposition}[Conditional multi-source reduction]\label{prop:multisource-reduction}
Fix $k$ and assume
\begin{enumerate}[label=(\alph*),nosep]
    \item $n_k\geq3$ and the induced subgraph $\G[\V_k]$ is connected;
    \item the team gain condition
    $\gamma_k\coloneqq\alpha_k-2\beta_kn_kf_{D_{\max},k}/R_k>0$;
    \item team-$k$ measurements are of the form $\sigma_i=f(p_i)+\eta_i$ with
    $\abs{\eta_i}\leq\delta$ and $\abs{\sigma_i}\leq f_{D_{\max},k}$;
    \item the certificate \eqref{eq:basin-certificate} holds for all $t\geq t_s$.
\end{enumerate}
Then team $k$ is an autonomous copy of the baseline system with parameters
$(n_k,R_k,\alpha_k,\beta_k,p_s^{(k)})$, and
\begin{equation}\label{eq:team-ultimate-bound}
    \limsup_{t\to\infty}\norm{p_k^*(t)-p_s^{(k)}}\leq\frac{\delta}{\kappa R_k},
\end{equation}
or $\varepsilon_k$ of \eqref{eq:team-epsilon-scale} if only a fraction $\chi_k$ is informed.
\end{proposition}

\begin{proof}
\emph{Step 1 (localise the robots).} By (d) and \cref{lem:basin-certificate}, every
$p_i(t)$ with $i\in\V_k$ lies in $\operatorname{int}\mathcal{B}_k$ for $t\geq t_s$.

\emph{Step 2 (the envelope collapses to one bowl).} Let $p\in\operatorname{int}\mathcal{B}_k$, so
$\norm{p-p_s^{(k)}}<\norm{p-p_s^{(j)}}$ for all $j\neq k$. Since $\kappa>0$ is the same for all
sources and $t\mapsto\kappa t^2$ is strictly increasing on $[0,\infty)$,
\begin{equation*}
    f_k(p)=\kappa\norm{p-p_s^{(k)}}^2<\kappa\norm{p-p_s^{(j)}}^2=f_j(p)
    \quad\forall j\neq k,
\end{equation*}
so the minimum in \eqref{eq:multisource-field} is attained uniquely at $k$ and
$f(p)=f_k(p)$. Equality of curvature is used here and only here; with source-dependent
$\kappa_j$ the ordering of squared distances no longer implies the ordering of field values and
the step fails.

\emph{Step 3 (measurements become single-source measurements).} Combining Steps 1 and 2, for
$i\in\V_k$ and $t\geq t_s$,
\begin{equation*}
    \sigma_i=f(p_i)+\eta_i=f_k(p_i)+\eta_i=\kappa\norm{p_i-p_s^{(k)}}^2+\eta_i,
\end{equation*}
which is exactly the measurement model of \cref{fact:localisation} for a single source at
$p_s^{(k)}$.

\emph{Step 4 (the dynamics decouple).} Robot $i\in\V_k$ uses
$\mathcal{N}_i\cap\V_k$ in \eqref{eq:baseline-controller}, so the consensus term involves only
team-$k$ states, and by Step 3 the measurement term involves only $p_i$ and $p_s^{(k)}$. Hence
$\dot p_i$ for $i\in\V_k$ is a function of $\{p_j\}_{j\in\V_k}$ alone: the team-$k$ closed loop
is autonomous and identical in form to \eqref{eq:baseline-controller}.

\emph{Step 5 (formation).} Condition (a) supplies connectivity and $n_k\geq3$, condition (c)
supplies the uniform measurement bound, and condition (b) is \eqref{eq:gain-condition} for the
team parameters. \cref{fact:formation} therefore applies to team $k$ and gives exact formation
after a finite time $T_{f,k}\leq\sqrt2\,n_k\sqrt{V_{f,k}(t_s)}/\gamma_k$.

\emph{Step 6 (localisation).} After $T_{f,k}$ the hypotheses of \cref{fact:localisation} hold
with $(n_k,R_k,\beta_k,p_s^{(k)})$, giving
$\dot e_k=-2\beta_k\kappa e_k+d_{\eta,k}$ with
$\norm{d_{\eta,k}}\leq2\beta_k\delta/R_k$ and therefore
$\limsup_t\norm{e_k}\leq\delta/(\kappa R_k)$, which is \eqref{eq:team-ultimate-bound}. Under
partial information, \cref{fact:epsilon} and \cref{lem:team-scaling} replace this by
$\varepsilon_k$.
\end{proof}

\begin{remark}[Why the proposition is conditional, and what is missing]\label{rem:gap1-conditional}
Hypothesis (d) is an assumption about the closed loop, not a consequence of it. The difficulty is
visible in the order of the proof: Step 5 establishes that $q_k(t)\to0$ in finite time, but
\eqref{eq:basin-certificate} needs $q_k$ to be small \emph{already}, over
$[t_s,T_{f,k}]$, in order to license Step 3. The two are not circular only because (d) is
assumed. Closing this loop requires either a certified bound on the post-split formation
transient $q_k$ on $[t_s,T_{f,k}]$, or a controller that keeps robots inside their basin during
the transient by construction. Neither is provided by the baseline controller, so
\cref{prop:multisource-reduction} should be read as a decomposition theorem with a run-time
checkable hypothesis, and \eqref{eq:basin-certificate} should be logged for every run rather than
assumed.
\end{remark}

\subsection{Numerical verification for Gap 1}\label{sec:gap1-numerics}

\paragraph{Check 1: cancellation of $n_k$.}
$\varepsilon_k$ was evaluated from the unsimplified expression of Step 1 of
\cref{lem:team-scaling}, for $n_k=3,\ldots,31$ at four fixed fractions
$\chi_k\in\{0.5,0.7,0.9,1.0\}$, with $\delta=0.05$, $\kappa=1$, $R_k=1$. Since
\eqref{eq:team-epsilon-scale} predicts no $n_k$ dependence, the test statistic is the relative
spread across $n_k$ at fixed $\chi_k$.

\begin{center}
\begin{longtable}{@{}p{0.34\linewidth}p{0.26\linewidth}p{0.30\linewidth}@{}}
\toprule
Quantity & Predicted & Measured \\
\midrule
\endhead
Relative spread of $\varepsilon_k$ over $n_k$ & $0$ (exact cancellation) & $5.6\times10^{-16}$ \\
$\varepsilon_k$ at $\chi_k=1$ & $\delta/(\kappa R_k)=0.05$ & $0.05$ \\
\bottomrule
\end{longtable}
\end{center}

The spread is at the level of double-precision rounding, so the cancellation in Steps 3--4 is
exact rather than approximate. The practical reading is the one in
\cref{rem:team-sizing}: a $12$-robot, $3$-source split cannot beat a $9$-robot, $3$-source split
on asymptotic accuracy at equal $(\chi_k,R_k)$; any observed advantage must come from transients
or from robustness, not from \eqref{eq:team-epsilon-scale}.

\paragraph{Check 2: positivity of the denominator.}
The factor $2\pi n_{\mathcal X}-n\abs{\sin(2\pi n_{\mathcal X}/n)}$ of \cref{lem:denominator} was
swept over all integer pairs with $3\leq n\leq60$, $1\leq n_{\mathcal X}\leq n$. The minimum is
$1.148\times10^{-2}$, attained at $(n,n_{\mathcal X})=(60,1)$, and no sign change occurs anywhere
on the grid. Separately,
$\max_{x\in[0,2\pi]}\left(\abs{\sin x}-x\right)=0$ to machine precision, confirming the
inequality used in the proof. The worst case is a single informed robot in a large team, which is
also where the bound is least useful: $\varepsilon$ is then large, but it is finite.

\paragraph{Check 3: the containment certificate.}
Because \eqref{eq:basin-certificate} is sufficient and not necessary
(\cref{rem:basin-sufficiency}), a single passing run proves nothing. Two two-source
configurations were run with $R=1$, and both the certificate margin
$d_k-\left(\norm{p_k^*-p_s^{(k)}}+R_k+q_k\right)$ and the ground-truth Voronoi margin
$m_k\coloneqq\min_{i\in\V_k}\min_{j\neq k}\left(\norm{p_i-p_s^{(j)}}-\norm{p_i-p_s^{(k)}}\right)$
were logged, the latter being positive exactly when containment holds.

\begin{center}
\begin{longtable}{@{}p{0.24\linewidth}p{0.16\linewidth}p{0.20\linewidth}p{0.16\linewidth}p{0.16\linewidth}@{}}
\toprule
Configuration & $d_k$ & Certificate & $\min_tm_k$ & Contained \\
\midrule
\endhead
Sources at $x=\pm4$ & $4.00$ & holds & $+4.90$ & yes \\
Sources at $x=\pm1.35$ & $1.35$ & fails & $-0.049$ & no \\
\bottomrule
\end{longtable}
\end{center}

In the separated case the certificate holds and containment holds, as the lemma requires. In the
tight case the ring radius $R=1$ alone consumes most of the clearance $d_k=1.35$, the certificate
fails, and containment is in fact lost during the transient. The lemma is never violated in the
direction it claims, and the second row shows that seam crossing is a real failure mode rather
than a theoretical worry. This is the empirical basis for the recommendation in
\cref{rem:gap1-conditional} to log \eqref{eq:basin-certificate} per run.

\section{Gap 2: a moving source}\label{sec:gap2}
\NotationContext{moving}

\subsection{Setting}\label{sec:gap2-setting}

Let the source follow an absolutely continuous trajectory $p_s(t)$ with
\begin{equation}\label{eq:speed-bound}
    \norm{\dot p_s(t)}\leq v_{\max}\quad\text{for almost every }t,
\end{equation}
and let the field be the moving bowl $f(p,t)=\kappa\norm{p-p_s(t)}^2$. Measurements are
$\sigma_i=f(p_i,t)+\eta_i$ with $\abs{\eta_i}\leq\delta$, and all robots are informed unless
stated otherwise. Write $e(t)=p^*(t)-p_s(t)$ and $\lambda=2\beta\kappa$.

\subsection{Step-by-step: source motion does not affect formation}\label{sec:gap2-formation}

The first thing to establish is that \cref{fact:formation} survives unchanged. This is not
obvious, because the measurement $\sigma_i$ now depends on time through $p_s(t)$.

\begin{lemma}[Formation is insensitive to source velocity]\label{lem:formation-invariance}
Assume $\G$ connected, $n\geq3$, and $\abs{\sigma_i(t)}\leq f_{D_{\max}}$ for all $i,t$. If
$\gamma=\alpha-2\beta nf_{D_{\max}}/R>0$, then $V_f(t)=0$ for all
$t\geq T_f\leq\sqrt2\,n\sqrt{V_f(0)}/\gamma$, and neither $T_f$ nor the bound depends on
$\dot p_s$.
\end{lemma}

\begin{proof}
\emph{Step 1 (shifted states move with the robots).} The slot offsets $R\phi_i$ are constant, so
$\dot z_i=\dot p_i=u_i$.

\emph{Step 2 (differentiate $V_f$).} With $\bar z=n^{-1}\sum_jz_j$ and $q_i=z_i-\bar z$,
\begin{equation*}
    \dot V_f=\sum_iq_i^{\top}\dot q_i
    =\sum_iq_i^{\top}\left(\dot z_i-\dot{\bar z}\right)
    =\sum_iq_i^{\top}\dot z_i-\Big(\underbrace{\sum_iq_i}_{=0}\Big)^{\!\top}\dot{\bar z}
    =\sum_iq_i^{\top}u_i.
\end{equation*}

\emph{Step 3 (insert the controller).} Substituting \eqref{eq:baseline-controller},
\begin{equation*}
    \dot V_f=\alpha\sum_iq_i^{\top}\sum_{j\in\mathcal{N}_i}\sgn\!\left(z_j-z_i\right)
    -\frac{2\beta}{R}\sum_i\sigma_i\,q_i^{\top}\phi_i.
\end{equation*}

\emph{Step 4 (the consensus term dissipates).} By the first identity of \cref{fact:sign} the
first sum equals $-S(z)$, so that term contributes $-\alpha S(z)$.

\emph{Step 5 (bound the measurement term).} Using $\abs{\sigma_i}\leq f_{D_{\max}}$,
$\norm{\phi_i}=1$, Cauchy--Schwarz, and then $Q(z)\leq nS(z)$ from \cref{fact:sign},
\begin{equation*}
    \abs*{\frac{2\beta}{R}\sum_i\sigma_iq_i^{\top}\phi_i}
    \leq\frac{2\beta f_{D_{\max}}}{R}\sum_i\norm{q_i}
    =\frac{2\beta f_{D_{\max}}}{R}Q(z)
    \leq\frac{2\beta nf_{D_{\max}}}{R}S(z).
\end{equation*}
This is the only place where $\sigma_i$ enters, and only its magnitude bound is used. Whether the
bowl is stationary or moving is irrelevant, which is the substance of the lemma.

\emph{Step 6 (combine).} Steps 4 and 5 give
\begin{equation}\label{eq:vf-decrease}
    \dot V_f\leq-\left(\alpha-\frac{2\beta nf_{D_{\max}}}{R}\right)S(z)=-\gamma S(z).
\end{equation}

\emph{Step 7 (finite-time convergence).} By the third identity of \cref{fact:sign},
$S(z)\geq\sqrt{2V_f}/n$, so $\dot V_f\leq-\gamma\sqrt{2V_f}/n$. Where $V_f>0$, put
$W\coloneqq\sqrt{V_f}$; then
\begin{equation*}
    \dot W=\frac{\dot V_f}{2\sqrt{V_f}}
    \leq-\frac{\gamma\sqrt2\sqrt{V_f}}{2n\sqrt{V_f}}
    =-\frac{\gamma}{\sqrt2\,n}.
\end{equation*}
Integrating, $W(t)\leq W(0)-\gamma t/(\sqrt2\,n)$, so $W$ reaches zero no later than
$T_f=\sqrt2\,n\,W(0)/\gamma=\sqrt2\,n\sqrt{V_f(0)}/\gamma$, and $V_f\equiv0$ afterwards because
$\dot V_f\leq0$. No step involved $\dot p_s$.
\end{proof}

\subsection{Step-by-step: the centroid error dynamics}\label{sec:gap2-dynamics}

\begin{lemma}[Error dynamics under source motion]\label{lem:moving-error-dynamics}
For $t\geq T_f$, with all robots informed,
\begin{equation}\label{eq:moving-error-dynamics}
    \dot e=-\lambda e+d_\eta-\dot p_s,
    \qquad
    \lambda=2\beta\kappa,
    \qquad
    d_\eta=-\frac{2\beta}{nR}\sum_i\eta_i\phi_i,
    \qquad
    \norm{d_\eta}\leq\frac{2\beta\delta}{R}.
\end{equation}
\end{lemma}

\begin{proof}
\emph{Step 1 (average the controller).} Since $p^*=n^{-1}\sum_ip_i$,
\begin{equation*}
    \dot p^*=\frac1n\sum_iu_i
    =\frac{\alpha}{n}\sum_i\sum_{j\in\mathcal{N}_i}\sgn\!\left(z_j-z_i\right)
    -\frac{2\beta}{nR}\sum_i\sigma_i\phi_i.
\end{equation*}

\emph{Step 2 (the consensus term averages to zero).} The double sum ranges over ordered pairs
$(i,j)$ with $\{i,j\}\in\E$. Each undirected edge contributes both
$\sgn(z_j-z_i)$ and $\sgn(z_i-z_j)$, and $\sgn$ is odd componentwise, so the two cancel. Hence
the first term vanishes exactly, for any graph and any states.

\emph{Step 3 (exact formation).} For $t\geq T_f$, \cref{lem:formation-invariance} gives $q_i=0$,
i.e.\ $z_i=\bar z$ for all $i$, which means $p_i=p^*+R\phi_i$.

\emph{Step 4 (expand the measurement).} Write $p_i-p_s=e+R\phi_i$. Then
\begin{align*}
    f(p_i,t)&=\kappa\norm{e+R\phi_i}^2
    =\kappa\left(\norm{e}^2+2Re^{\top}\phi_i+R^2\norm{\phi_i}^2\right)\\
    &=\kappa\norm{e}^2+2\kappa Re^{\top}\phi_i+\kappa R^2 .
\end{align*}

\emph{Step 5 (project onto the slot directions).} Multiply by $\phi_i$ and sum, using
\cref{fact:moments}:
\begin{equation*}
    \sum_if(p_i,t)\phi_i
    =\kappa\norm{e}^2\underbrace{\sum_i\phi_i}_{=0}
    +2\kappa R\underbrace{\left(\sum_i\phi_i\phi_i^{\top}\right)}_{=\frac n2I_2}e
    +\kappa R^2\underbrace{\sum_i\phi_i}_{=0}
    =\kappa nRe.
\end{equation*}
The two constant-in-$i$ terms are annihilated by the first moment identity; only the term linear
in $e$ survives, and it survives \emph{exactly}, with no truncation. This is why the quadratic
bowl gives an unbiased centroid feedback.

\emph{Step 6 (add noise).} With $\sigma_i=f(p_i,t)+\eta_i$,
$\sum_i\sigma_i\phi_i=\kappa nRe+\sum_i\eta_i\phi_i$, so from Step 1
\begin{equation*}
    \dot p^*=-\frac{2\beta}{nR}\left(\kappa nRe+\sum_i\eta_i\phi_i\right)
    =-2\beta\kappa e+d_\eta .
\end{equation*}

\emph{Step 7 (differentiate the error).} $e=p^*-p_s$ gives
$\dot e=\dot p^*-\dot p_s=-\lambda e+d_\eta-\dot p_s$.

\emph{Step 8 (bound the noise term).} By the triangle inequality,
$\norm{d_\eta}\leq\frac{2\beta}{nR}\sum_i\abs{\eta_i}\norm{\phi_i}\leq\frac{2\beta}{nR}\cdot n\delta
=\frac{2\beta\delta}{R}$.
\end{proof}

\begin{remark}[The structural point]\label{rem:moving-structure}
Comparing \eqref{eq:moving-error-dynamics} with \eqref{eq:baseline-error-dynamics}, source motion
enters the error dynamics as an additive disturbance $-\dot p_s$ on an otherwise unchanged
exponentially stable linear system, and it does not enter the formation subsystem at all
(\cref{lem:formation-invariance}). This is the whole reason Gap 2 admits a clean answer: the
baseline is input-to-state stable with respect to that channel, so no redesign is needed to get a
bound, only to get a \emph{small} bound.
\end{remark}

\subsection{Step-by-step: the ISS bound}\label{sec:gap2-iss}

\begin{theorem}[Ultimate bound under bounded source motion]\label{thm:moving-iss}
Under \eqref{eq:speed-bound} and the hypotheses of \cref{lem:moving-error-dynamics}, for all
$t\geq T_f$
\begin{equation}\label{eq:moving-transient}
    \norm{e(t)}\leq e^{-\lambda(t-T_f)}\norm{e(T_f)}
    +\left(\frac{\delta}{\kappa R}+\frac{v_{\max}}{2\beta\kappa}\right)
    \left(1-e^{-\lambda(t-T_f)}\right),
\end{equation}
and therefore
\begin{equation}\label{eq:moving-ultimate-bound}
    \limsup_{t\to\infty}\norm{e(t)}
    \leq\frac{\delta}{\kappa R}+\frac{v_{\max}}{2\beta\kappa}.
\end{equation}
\end{theorem}

\begin{proof}
\emph{Step 1 (collect the disturbance).} Write \eqref{eq:moving-error-dynamics} as
$\dot e=-\lambda e+w$ with $w\coloneqq d_\eta-\dot p_s$. By \eqref{eq:speed-bound} and Step 8 of
\cref{lem:moving-error-dynamics},
\begin{equation}\label{eq:disturbance-bound}
    \norm{w(t)}\leq\frac{2\beta\delta}{R}+v_{\max}\eqqcolon c .
\end{equation}

\emph{Step 2 (a differential inequality for the norm).} Let $r(t)\coloneqq\norm{e(t)}$. Where
$r>0$, $r$ is differentiable with $\dot r=e^{\top}\dot e/\norm{e}$, so
\begin{equation*}
    \dot r=\frac{e^{\top}\left(-\lambda e+w\right)}{\norm{e}}
    =-\lambda r+\frac{e^{\top}w}{\norm{e}}
    \leq-\lambda r+\norm{w}
    \leq-\lambda r+c,
\end{equation*}
using Cauchy--Schwarz and \eqref{eq:disturbance-bound}.

\emph{Step 3 (handle $r=0$).} At a time where $e=0$ the map $r$ need not be differentiable, but
its upper right Dini derivative satisfies
$D^+r\leq\norm{\dot e}=\norm{w}\leq c=-\lambda\cdot0+c$. Hence
\begin{equation}\label{eq:dini}
    D^+r(t)\leq-\lambda r(t)+c
\end{equation}
holds for all $t\geq T_f$ without exception, and $r$ is locally Lipschitz, so the comparison
principle applies to \eqref{eq:dini}.

\emph{Step 4 (integrate the comparison system).} Let $\bar r$ solve
$\dot{\bar r}=-\lambda\bar r+c$, $\bar r(T_f)=r(T_f)$. Multiplying by $e^{\lambda t}$,
$\frac{d}{dt}\left(e^{\lambda t}\bar r\right)=ce^{\lambda t}$; integrating from $T_f$ to $t$,
\begin{equation*}
    e^{\lambda t}\bar r(t)-e^{\lambda T_f}r(T_f)
    =\frac{c}{\lambda}\left(e^{\lambda t}-e^{\lambda T_f}\right),
\end{equation*}
and multiplying by $e^{-\lambda t}$,
\begin{equation}\label{eq:comparison-solution}
    \bar r(t)=e^{-\lambda(t-T_f)}r(T_f)
    +\frac{c}{\lambda}\left(1-e^{-\lambda(t-T_f)}\right).
\end{equation}
By comparison, $r(t)\leq\bar r(t)$ for $t\geq T_f$.

\emph{Step 5 (evaluate $c/\lambda$).} This is where the two disturbance sources separate:
\begin{equation}\label{eq:c-over-lambda}
    \frac{c}{\lambda}
    =\frac{2\beta\delta/R+v_{\max}}{2\beta\kappa}
    =\frac{2\beta\delta}{2\beta\kappa R}+\frac{v_{\max}}{2\beta\kappa}
    =\frac{\delta}{\kappa R}+\frac{v_{\max}}{2\beta\kappa}.
\end{equation}
The first term is exactly the static radius of \cref{fact:localisation}; the second is the price
of motion. Note that $\beta$ cancels from the noise term but not from the motion term, so
increasing $\beta$ shrinks the motion penalty without touching the noise floor.

\emph{Step 6 (conclude).} Substituting \eqref{eq:c-over-lambda} into
\eqref{eq:comparison-solution} gives \eqref{eq:moving-transient}, and letting $t\to\infty$ with
$e^{-\lambda(t-T_f)}\to0$ gives \eqref{eq:moving-ultimate-bound}.
\end{proof}

\subsection{Step-by-step: the exact constant-velocity lag}\label{sec:gap2-lag}

\eqref{eq:moving-ultimate-bound} is an inequality on a norm, so it cannot be falsified by a
single trajectory. The following corollary is an equality on a vector, and is therefore the
sharper test.

\begin{corollary}[Constant-velocity steady lag]\label{cor:constant-velocity}
Let $\dot p_s\equiv v$ be constant and $\eta_i\equiv0$. Then for $t\geq T_f$
\begin{equation}\label{eq:constant-velocity-solution}
    e(t)=e^{-\lambda(t-T_f)}e(T_f)
    -\frac{v}{\lambda}\left(1-e^{-\lambda(t-T_f)}\right),
\end{equation}
so the error converges to the constant vector
\begin{equation}\label{eq:constant-velocity-lag}
    e_\infty=-\frac{v}{\lambda}=-\frac{v}{2\beta\kappa}.
\end{equation}
\end{corollary}

\begin{proof}
\emph{Step 1 (specialise the dynamics).} With $\eta_i\equiv0$, Step 6 of
\cref{lem:moving-error-dynamics} gives $d_\eta=0$, so \eqref{eq:moving-error-dynamics} becomes
the linear time-invariant system $\dot e=-\lambda e-v$ with constant input.

\emph{Step 2 (integrating factor).} Because $\lambda$ is constant,
\begin{equation*}
    \frac{d}{dt}\left(e^{\lambda t}e(t)\right)
    =e^{\lambda t}\left(\dot e+\lambda e\right)
    =-e^{\lambda t}v .
\end{equation*}

\emph{Step 3 (integrate componentwise).} Integrating from $T_f$ to $t$, and using that $v$ is
constant,
\begin{equation*}
    e^{\lambda t}e(t)-e^{\lambda T_f}e(T_f)
    =-v\int_{T_f}^te^{\lambda s}\,ds
    =-\frac{v}{\lambda}\left(e^{\lambda t}-e^{\lambda T_f}\right).
\end{equation*}

\emph{Step 4 (solve).} Multiplying by $e^{-\lambda t}$ yields
\eqref{eq:constant-velocity-solution}.

\emph{Step 5 (limit).} As $t\to\infty$ the factor $e^{-\lambda(t-T_f)}$ vanishes since
$\lambda=2\beta\kappa>0$, leaving $e(t)\to-v/\lambda$, which is
\eqref{eq:constant-velocity-lag}.

\emph{Step 6 (consistency).} Taking norms,
$\norm{e_\infty}=\norm{v}/(2\beta\kappa)\leq v_{\max}/(2\beta\kappa)$, so
\eqref{eq:constant-velocity-lag} saturates the motion term of
\eqref{eq:moving-ultimate-bound} while contributing nothing to the noise term, as expected for
$\delta=0$. The bound of \cref{thm:moving-iss} is therefore tight in this regime.
\end{proof}

\begin{remark}[Direction, not just magnitude]\label{rem:lag-direction}
\eqref{eq:constant-velocity-lag} states that the formation trails the source along
$-v$ exactly, with gain $1/(2\beta\kappa)$. It is a two-component prediction, so it fixes both
the magnitude and the direction of the steady error. Any implementation error that rotates the
effective feedback, such as a slot-assignment mismatch or an inconsistent $\phi_i$ ordering, will
break the direction while possibly preserving the magnitude, which makes this the most
informative single test for Gap 2.
\end{remark}

\subsection{Partial information and velocity feedforward}\label{sec:gap2-partial}

Two extensions follow with no new machinery.

\paragraph{Partial information.}
If only a guaranteed fraction of robots is informed, Step 5 of
\cref{lem:moving-error-dynamics} no longer gives the exact identity $\sum_i\sigma_i\phi_i=\kappa nRe$
and \cref{fact:epsilon} must be used in place of \cref{fact:localisation}. The comparison
argument of \cref{thm:moving-iss} is unchanged in structure, and yields
\begin{equation}\label{eq:moving-partial-bound}
    \limsup_{t\to\infty}\norm{e(t)}
    \leq\varepsilon\!\left(\underline n_{\mathcal X}\right)+\frac{v_{\max}}{\lambda_{\mathcal X}},
\end{equation}
where $\underline n_{\mathcal X}$ is a lower bound on the informed count and
$\lambda_{\mathcal X}$ is the corresponding contraction rate. This is conditional in a way
\eqref{eq:moving-ultimate-bound} is not: source motion changes which robots are within sensing
range, so $\underline n_{\mathcal X}$ is itself a closed-loop quantity and must be verified, not
assumed.

\paragraph{Feedforward.}
Suppose an estimate $\hat v(t)$ of $\dot p_s$ is broadcast and every robot applies
$u_i^{\mathrm{ff}}=u_i+\hat v(t)$. Two consequences follow immediately.
\begin{enumerate}[label=(\roman*),nosep]
    \item \emph{Formation is untouched.} The added term is common to all robots, so it cancels
    from $\dot z_i-\dot z_j$ and hence from $\dot q_i$. \cref{lem:formation-invariance} applies
    verbatim.
    \item \emph{The disturbance shrinks.} Averaging as in Step 1 of
    \cref{lem:moving-error-dynamics} adds $\hat v$ to $\dot p^*$, so
    $\dot e=-\lambda e+d_\eta+\tilde v$ with $\tilde v\coloneqq\hat v-\dot p_s$. If
    $\norm{\tilde v}\leq\bar v$, then \cref{thm:moving-iss} applies with $v_{\max}$ replaced by
    $\bar v$, and perfect feedforward recovers the static radius $\delta/(\kappa R)$ exactly.
\end{enumerate}
Supplying $\hat v$ by a distributed estimator is a separate design problem and is not solved by
the baseline controller; time-varying gradient-free cooperative seeking under different
structural assumptions is treated in \cite{michael2023}.

\subsection{Numerical verification for Gap 2}\label{sec:gap2-numerics}

\paragraph{Check 4: the vector lag.}
A noiseless run with $n=6$, $R=1$, $\kappa=1$, $\beta=0.05$ (so $\lambda=0.1$) and
$v=(0.08,0.04)\,\mathrm{m/s}$ was integrated to steady state, and both components of $e$ were
compared with \eqref{eq:constant-velocity-lag}.

\begin{center}
\begin{longtable}{@{}p{0.30\linewidth}p{0.30\linewidth}p{0.30\linewidth}@{}}
\toprule
Quantity & Predicted & Measured \\
\midrule
\endhead
$e_\infty=-v/(2\beta\kappa)$ & $(-0.800,\,-0.400)$ & $(-0.818,\,-0.410)$ \\
Relative error on the vector & $0$ & $2.31\%$ \\
Slope of $\norm{e_\infty}$ vs $\norm{v}$ & $1/\lambda=10$ & within $4.97\%$ \\
\bottomrule
\end{longtable}
\end{center}

Both components match, so the direction is confirmed and not merely the magnitude, which is the
test \cref{rem:lag-direction} asks for. The residual few percent is the fixed-step integration
error of the sign controller, which chatters about the consensus manifold at an amplitude
proportional to the step size; the speeds above were chosen so the predicted lag is well clear of
that floor. At smaller speeds the predicted lag and the numerical floor are comparable and the
experiment cannot resolve \eqref{eq:constant-velocity-lag}, which is a property of the
discretisation rather than of \cref{cor:constant-velocity}.

\paragraph{Check 5: the affine and inverse dependencies.}
Sweeping $\norm{v}$ at fixed $\beta\kappa$ reproduces the predicted slope $1/\lambda$ to within
$5\%$, confirming that the steady lag is affine in speed and inverse in $\beta\kappa$, which is
the design trade-off identified in Step 5 of \cref{thm:moving-iss}.

\section{Gap 3: escaping local extrema and certifying the global maximum}\label{sec:gap3}
\NotationContext{field}

\subsection{Setting}\label{sec:gap3-setting}

Let $\Omega\subset\R^2$ be compact and $f\in C^3(\Omega)$ be a nonconvex field to be maximised,
with isolated local maximisers $\ell_1,\ldots,\ell_P$ ordered so that
$f(\ell_1)>f(\ell_2)\geq\cdots\geq f(\ell_P)$, global value $f^\star=f(\ell_1)$, and value gap
\begin{equation}\label{eq:value-gap}
    \Delta\coloneqq f(\ell_1)-f(\ell_2)>0.
\end{equation}
Let $L_f$ be a Lipschitz constant of $f$ on $\Omega$ and $L_H$ one of $\nabla^2f$, so
\begin{equation}\label{eq:hessian-lipschitz}
    \norm{\nabla^2f(x)-\nabla^2f(y)}\leq L_H\norm{x-y}
    \quad\forall x,y\in\Omega .
\end{equation}
For maximisation the measurement term in \eqref{eq:baseline-controller} carries the opposite
sign, and the ring of radius $R$ around the centroid supplies the circular gradient estimate
\begin{equation}\label{eq:circular-estimator}
    g_R(p)\coloneqq\frac{2}{nR}\sum_{i=1}^nf\!\left(p+R\phi_i\right)\phi_i .
\end{equation}

\subsection{Step-by-step: the estimator bias is $O(R^2)$, but only for $n\geq4$}\label{sec:gap3-bias}

\begin{lemma}[Circular gradient estimator bias]\label{lem:estimator-bias}
Let $f\in C^3$ satisfy \eqref{eq:hessian-lipschitz}, let $n\geq4$, and let
$B(p,R)\subset\Omega$. Then
\begin{equation}\label{eq:estimator-bias}
    g_R(p)=\nabla f(p)+r_R(p),
    \qquad
    \norm{r_R(p)}\leq\frac{L_HR^2}{3}.
\end{equation}
\end{lemma}

\begin{proof}
\emph{Step 1 (Taylor with integral remainder).} Fix $h\in\R^2$ with $p+sh\in\Omega$ for
$s\in[0,1]$. The exact second-order Taylor formula is
\begin{equation}\label{eq:taylor-exact}
    f(p+h)=f(p)+\nabla f(p)^{\top}h
    +\int_0^1(1-s)\,h^{\top}\nabla^2f(p+sh)\,h\,ds .
\end{equation}

\emph{Step 2 (split off the frozen Hessian).} Since $\int_0^1(1-s)\,ds=\frac12$, subtracting and
adding $\frac12h^{\top}\nabla^2f(p)h$ inside \eqref{eq:taylor-exact} gives
\begin{equation}\label{eq:taylor-split}
    f(p+h)=f(p)+\nabla f(p)^{\top}h+\frac12h^{\top}\nabla^2f(p)h+\rho(h),
\end{equation}
\begin{equation}\label{eq:remainder-def}
    \rho(h)=\int_0^1(1-s)\,h^{\top}
    \left[\nabla^2f(p+sh)-\nabla^2f(p)\right]h\,ds .
\end{equation}

\emph{Step 3 (bound the remainder).} Applying \eqref{eq:hessian-lipschitz} to the bracket with
$x=p+sh$, $y=p$, so that $\norm{x-y}=s\norm{h}$,
\begin{equation*}
    \abs{\rho(h)}
    \leq\int_0^1(1-s)\,L_Hs\norm{h}\cdot\norm{h}^2\,ds
    =L_H\norm{h}^3\int_0^1s(1-s)\,ds
    =\frac{L_H\norm{h}^3}{6},
\end{equation*}
using $\int_0^1s(1-s)\,ds=\frac12-\frac13=\frac16$.

\emph{Step 4 (insert into the estimator).} Put $h=R\phi_i$, so $\norm{h}=R$, and write
$\rho_i\coloneqq\rho(R\phi_i)$. Substituting \eqref{eq:taylor-split} into
\eqref{eq:circular-estimator} and separating the four terms,
\begin{align*}
    g_R(p)=\frac{2}{nR}\Bigg[
    &f(p)\sum_i\phi_i
    +R\sum_i\left(\nabla f(p)^{\top}\phi_i\right)\phi_i\\
    &+\frac{R^2}{2}\sum_i\left(\phi_i^{\top}\nabla^2f(p)\phi_i\right)\phi_i
    +\sum_i\rho_i\phi_i\Bigg].
\end{align*}

\emph{Step 5 (kill the zeroth-order term).} $\sum_i\phi_i=0$ by \cref{fact:moments}, so the
first bracketed term vanishes. This is what makes the estimator insensitive to an unknown
additive offset in the field.

\emph{Step 6 (recover the gradient exactly).} Since
$\left(\nabla f^{\top}\phi_i\right)\phi_i=\phi_i\phi_i^{\top}\nabla f$,
\begin{equation*}
    \frac{2}{nR}\cdot R\sum_i\phi_i\phi_i^{\top}\nabla f(p)
    =\frac{2}{n}\cdot\frac n2I_2\,\nabla f(p)
    =\nabla f(p),
\end{equation*}
using the second identity of \cref{fact:moments}. The gradient is reproduced with no constant to
tune.

\emph{Step 7 (kill the curvature term).} Here $n\geq4$ is used: by \eqref{eq:third-moment} with
$H=\nabla^2f(p)$ symmetric, $\sum_i\left(\phi_i^{\top}H\phi_i\right)\phi_i=0$, so the third term
vanishes identically. Note that it is the \emph{third} moment of the slot directions that is
required, not the second, which is why $n=3$ is excluded.

\emph{Step 8 (bound what is left).} The only surviving term is
$r_R(p)=\frac{2}{nR}\sum_i\rho_i\phi_i$, and by Step 3 with $\norm{h}=R$,
\begin{equation*}
    \norm{r_R(p)}
    \leq\frac{2}{nR}\sum_i\abs{\rho_i}\norm{\phi_i}
    \leq\frac{2}{nR}\cdot n\cdot\frac{L_HR^3}{6}
    =\frac{L_HR^2}{3}. \qedhere
\end{equation*}
\end{proof}

The exclusion of $n=3$ is not an artefact of the proof. The next lemma computes the leading term
that survives, which turns the $n=3$ case from a gap in the argument into a quantitative
prediction.

\begin{lemma}[Exact leading bias at $n=3$]\label{lem:n3-bias}
Let $n=3$ and write $H\coloneqq\nabla^2f(p)=\begin{psmallmatrix}H_{11}&H_{12}\\H_{12}&H_{22}\end{psmallmatrix}$.
Then
\begin{equation}\label{eq:n3-bias-vector}
    r_R(p)=\frac R4
    \begin{pmatrix}H_{11}-H_{22}\\-2H_{12}\end{pmatrix}
    +O\!\left(R^2\right),
    \qquad
    \norm{r_R(p)}=\frac R4\sqrt{\left(H_{11}-H_{22}\right)^2+4H_{12}^2}+O\!\left(R^2\right).
\end{equation}
In particular the bias is $O(R)$ rather than $O(R^2)$, and it vanishes to first order exactly
when $H\propto I_2$.
\end{lemma}

\begin{proof}
\emph{Step 1 (reduce the quadratic form).} With $\phi=(\cos\theta,\sin\theta)$ and the
double-angle identities $\cos^2\theta=\frac{1+\cos2\theta}{2}$,
$\sin^2\theta=\frac{1-\cos2\theta}{2}$, $2\cos\theta\sin\theta=\sin2\theta$,
\begin{equation*}
    \phi^{\top}H\phi
    =H_{11}\cos^2\theta+2H_{12}\cos\theta\sin\theta+H_{22}\sin^2\theta
    =\frac{H_{11}+H_{22}}{2}
    +\frac{H_{11}-H_{22}}{2}\cos2\theta
    +H_{12}\sin2\theta .
\end{equation*}

\emph{Step 2 (the isotropic part drops out).} The first term is independent of $i$, so it
multiplies $\sum_i\phi_i=0$ and contributes nothing.

\emph{Step 3 (evaluate the $\cos2\theta$ moment).} For $\theta_i\in\{0,2\pi/3,4\pi/3\}$ we have
$3\theta_i\in\{0,2\pi,4\pi\}$, hence $\sum_i\cos3\theta_i=3$ and $\sum_i\sin3\theta_i=0$, while
$\sum_i\cos\theta_i=\sum_i\sin\theta_i=0$. Using
$\cos2\theta\cos\theta=\frac12(\cos3\theta+\cos\theta)$ and
$\cos2\theta\sin\theta=\frac12(\sin3\theta-\sin\theta)$,
\begin{equation*}
    \sum_i\cos2\theta_i\,\phi_i
    =\left(\tfrac12(3+0),\ \tfrac12(0-0)\right)
    =\left(\tfrac32,\,0\right).
\end{equation*}

\emph{Step 4 (evaluate the $\sin2\theta$ moment).} Using
$\sin2\theta\cos\theta=\frac12(\sin3\theta+\sin\theta)$ and
$\sin2\theta\sin\theta=\frac12(\cos\theta-\cos3\theta)$,
\begin{equation*}
    \sum_i\sin2\theta_i\,\phi_i
    =\left(\tfrac12(0+0),\ \tfrac12(0-3)\right)
    =\left(0,\,-\tfrac32\right).
\end{equation*}

\emph{Step 5 (assemble the third moment).} Combining Steps 1--4,
\begin{equation}\label{eq:n3-third-moment}
    \sum_{i=1}^3\left(\phi_i^{\top}H\phi_i\right)\phi_i
    =\frac{H_{11}-H_{22}}{2}\begin{pmatrix}3/2\\0\end{pmatrix}
    +H_{12}\begin{pmatrix}0\\-3/2\end{pmatrix}
    =\frac34\begin{pmatrix}H_{11}-H_{22}\\-2H_{12}\end{pmatrix},
\end{equation}
which is nonzero unless $H_{11}=H_{22}$ and $H_{12}=0$. This is the promised failure of
\eqref{eq:third-moment} at $n=3$.

\emph{Step 6 (propagate through the estimator).} Steps 4--6 of \cref{lem:estimator-bias} are
unchanged, so the curvature term of Step 7 now survives:
\begin{equation*}
    r_R(p)=\frac{2}{nR}\cdot\frac{R^2}{2}\sum_i\left(\phi_i^{\top}H\phi_i\right)\phi_i
    +\frac{2}{nR}\sum_i\rho_i\phi_i
    =\frac R3\cdot\frac34\begin{pmatrix}H_{11}-H_{22}\\-2H_{12}\end{pmatrix}
    +O\!\left(R^2\right),
\end{equation*}
where $n=3$ was substituted and Step 8 of \cref{lem:estimator-bias} bounds the remainder by
$L_HR^2/3$. Simplifying $\frac R3\cdot\frac34=\frac R4$ gives the first claim.

\emph{Step 7 (take the norm).} $\norm{\left(a,-2b\right)}=\sqrt{a^2+4b^2}$ with
$a=H_{11}-H_{22}$, $b=H_{12}$, giving the second claim.
\end{proof}

\begin{remark}[Design consequence]\label{rem:n3-design}
\cref{lem:n3-bias} says a three-robot ring measures a systematic direction error whenever the
field is anisotropically curved, and that the error shrinks only linearly in $R$. Since $R$ also
sets the noise rejection through $\delta/(\kappa R)$, shrinking $R$ to fight the bias worsens the
noise floor. Using $n\geq4$ removes the trade-off at its source, which is a concrete reason to
prefer $n_k\geq4$ per team in \cref{sec:gap1} even though \cref{fact:formation} only needs
$n_k\geq3$.
\end{remark}

\subsection{Step-by-step: no controller of this form can be globally convergent}\label{sec:gap3-trapping}

\begin{proposition}[Local trapping is unavoidable and not exceptional]\label{prop:no-global-guarantee}
Consider the ideal ascent flow obtained from \eqref{eq:circular-estimator} in the limit
$R\to0$ with exact formation,
\begin{equation}\label{eq:ideal-ascent}
    \dot p^*=\beta\nabla f\!\left(p^*\right),
    \qquad\beta>0 .
\end{equation}
Let $\ell$ be a local maximiser with $\nabla^2f(\ell)\prec0$ and $f(\ell)<f^\star$. Then there is
an open set $\Omega_c\ni\ell$, of positive Lebesgue measure, that is forward invariant under
\eqref{eq:ideal-ascent} and from which every trajectory converges to $\ell$. Consequently no
choice of $\beta$ makes \eqref{eq:ideal-ascent} globally convergent to $\ell_1$.
\end{proposition}

\begin{proof}
\emph{Step 1 ($\ell$ is an equilibrium).} $\nabla f(\ell)=0$ because $\ell$ is an interior local
maximiser, so $\dot p^*=0$ there.

\emph{Step 2 ($\ell$ is locally exponentially stable).} The Jacobian of
\eqref{eq:ideal-ascent} at $\ell$ is $\beta\nabla^2f(\ell)$. As $\nabla^2f(\ell)$ is symmetric
and negative definite, its eigenvalues are real and negative, hence so are those of
$\beta\nabla^2f(\ell)$ for $\beta>0$. Local exponential stability follows.

\emph{Step 3 ($\ell$ is an isolated critical point).} $\nabla^2f(\ell)\prec0$ is nonsingular, so
by the inverse function theorem applied to $\nabla f$ there is a neighbourhood of $\ell$ on which
$\nabla f$ vanishes only at $\ell$.

\emph{Step 4 (a Lyapunov function).} Define $V(p)\coloneqq f(\ell)-f(p)$. Then $V(\ell)=0$ and,
shrinking to the neighbourhood of Step 3 where $\ell$ is the unique maximiser, $V(p)>0$ for
$p\neq\ell$. Along \eqref{eq:ideal-ascent},
\begin{equation}\label{eq:vdot-ascent}
    \dot V=-\nabla f(p^*)^{\top}\dot p^*
    =-\beta\norm{\nabla f(p^*)}^2\leq0 .
\end{equation}

\emph{Step 5 (construct the invariant set).} Choose $c>0$ small enough that the connected
component $\Omega_c$ of $\{p:V(p)<c\}$ containing $\ell$ is bounded, contained in the interior of
$\Omega$, and contains no critical point of $f$ other than $\ell$. Such a $c$ exists by Step 3
and continuity of $f$. By \eqref{eq:vdot-ascent}, $V$ is nonincreasing along trajectories, so
$V(p^*(t))<c$ persists and $\Omega_c$ is forward invariant.

\emph{Step 6 (convergence inside $\Omega_c$).} $\Omega_c$ is compact in closure and forward
invariant, so LaSalle's theorem applies: every trajectory starting in $\Omega_c$ converges to the
largest invariant subset of $\{p\in\Omega_c:\dot V=0\}$. By \eqref{eq:vdot-ascent} that set is
$\{p\in\Omega_c:\nabla f(p)=0\}=\{\ell\}$ by the choice of $c$. Hence $p^*(t)\to\ell$.

\emph{Step 7 ($\Omega_c$ is not negligible).} $V$ is continuous, so $\{V<c\}$ is open, and
$\Omega_c$ is an open connected component of an open set containing $\ell$. A nonempty open
subset of $\R^2$ has positive Lebesgue measure. The trapping set is therefore not a measure-zero
pathology that could be dismissed by a genericity argument.

\emph{Step 8 (conclude).} $f(\ell)<f^\star=f(\ell_1)$, so trajectories from $\Omega_c$ converge
to a point that is not the global maximiser, for every $\beta>0$. Since $\beta$ only rescales
time in \eqref{eq:ideal-ascent}, no gain tuning can help. Escaping $\Omega_c$ requires adding a
mechanism that is not a function of the local gradient alone.
\end{proof}

\begin{remark}[What this does and does not say]\label{rem:trapping-scope}
\cref{prop:no-global-guarantee} is a statement about \eqref{eq:ideal-ascent}, which is the
$R\to0$ ideal of the circular estimator. A finite ring radius $R$ averages the field over a
circle and can therefore step over narrow barriers, and \cref{lem:estimator-bias} quantifies how
far the finite-$R$ flow deviates from \eqref{eq:ideal-ascent}. So the proposition does not forbid
a finite-$R$ implementation from occasionally escaping; it forbids any \emph{guarantee} of global
convergence based on local ascent, and it makes the standard remedies, namely restarts and
dither, structural necessities rather than heuristics.
\end{remark}

\subsection{Step-by-step: multi-start selection identifies the global maximum}\label{sec:gap3-multistart}
\NotationContext{multistart}

Given \cref{prop:no-global-guarantee}, the practical question is not how to avoid traps but how
to \emph{detect} that the best of several outcomes is the global one. Run $M$ independent teams,
let $\hat\ell_k$ be the terminal centroid of team $k$, and let $\hat F_k$ be the score it
reports, for instance a formation-averaged field value.

\begin{theorem}[Correct selection under bounded errors]\label{thm:multistart}
Assume
\begin{enumerate}[label=(\alph*),nosep]
    \item at least one team $k_1$ terminates near the global maximiser, in the sense that
    $\norm{\hat\ell_{k_1}-\ell_1}\leq\varepsilon_{\mathrm{loc}}$;
    \item every other team terminates near \emph{some} local maximiser,
    $\norm{\hat\ell_k-\ell_{j(k)}}\leq\varepsilon_{\mathrm{loc}}$ for some index $j(k)$;
    \item scores are accurate to $\delta_F$, i.e.\ $\abs{\hat F_k-f(\hat\ell_k)}\leq\delta_F$ for
    all $k$;
    \item the separation condition
    \begin{equation}\label{eq:multistart-condition}
        \Delta>2\left(L_f\varepsilon_{\mathrm{loc}}+\delta_F\right).
    \end{equation}
\end{enumerate}
Then every maximiser of the reported scores is a team that terminated near $\ell_1$:
\begin{equation}\label{eq:multistart-selection}
    \argmax_{1\leq k\leq M}\hat F_k
    \subseteq\left\{k:\norm{\hat\ell_k-\ell_1}\leq\varepsilon_{\mathrm{loc}}\right\}.
\end{equation}
\end{theorem}

\begin{proof}
\emph{Step 1 (lower bound the winner's score).} By (c) applied to $k_1$,
$\hat F_{k_1}\geq f(\hat\ell_{k_1})-\delta_F$.

\emph{Step 2 (transfer to the true peak value).} $f$ is $L_f$-Lipschitz on $\Omega$, so by (a)
\begin{equation*}
    \abs{f(\hat\ell_{k_1})-f(\ell_1)}
    \leq L_f\norm{\hat\ell_{k_1}-\ell_1}
    \leq L_f\varepsilon_{\mathrm{loc}},
\end{equation*}
hence $f(\hat\ell_{k_1})\geq f(\ell_1)-L_f\varepsilon_{\mathrm{loc}}$.

\emph{Step 3 (combine).} Chaining Steps 1 and 2,
\begin{equation}\label{eq:multistart-lower}
    \hat F_{k_1}\geq f(\ell_1)-L_f\varepsilon_{\mathrm{loc}}-\delta_F .
\end{equation}

\emph{Step 4 (upper bound a suboptimal team).} Let $k$ be any team with $j(k)\geq2$. By (c) and
then (b) with the same Lipschitz argument,
\begin{equation*}
    \hat F_k\leq f(\hat\ell_k)+\delta_F
    \leq f\!\left(\ell_{j(k)}\right)+L_f\varepsilon_{\mathrm{loc}}+\delta_F .
\end{equation*}

\emph{Step 5 (use the value gap).} The ordering of the maximisers and
\eqref{eq:value-gap} give $f(\ell_{j})\leq f(\ell_2)=f(\ell_1)-\Delta$ for every $j\geq2$, so
\begin{equation}\label{eq:multistart-upper}
    \hat F_k\leq f(\ell_1)-\Delta+L_f\varepsilon_{\mathrm{loc}}+\delta_F .
\end{equation}

\emph{Step 6 (subtract).} Subtracting \eqref{eq:multistart-upper} from
\eqref{eq:multistart-lower}, the common term $f(\ell_1)$ cancels:
\begin{equation*}
    \hat F_{k_1}-\hat F_k
    \geq-L_f\varepsilon_{\mathrm{loc}}-\delta_F
    +\Delta-L_f\varepsilon_{\mathrm{loc}}-\delta_F
    =\Delta-2\left(L_f\varepsilon_{\mathrm{loc}}+\delta_F\right).
\end{equation*}

\emph{Step 7 (apply the separation condition).} By (d) the right-hand side is strictly positive,
so $\hat F_{k_1}>\hat F_k$ for every team $k$ with $j(k)\geq2$.

\emph{Step 8 (conclude).} No team that terminated near a suboptimal maximiser can attain the
maximum of $\{\hat F_k\}$, since $k_1$ strictly beats it. Every maximiser must therefore be a
team satisfying $\norm{\hat\ell_k-\ell_1}\leq\varepsilon_{\mathrm{loc}}$, which is
\eqref{eq:multistart-selection}.
\end{proof}

\begin{remark}[Reading the condition]\label{rem:multistart-reading}
\eqref{eq:multistart-condition} is the worst case in which both error sources conspire against
the winner and in favour of the runner-up simultaneously, each at full magnitude and in the least
favourable direction. It is sufficient, not necessary, and in practice selection succeeds well
outside it, as \cref{sec:gap3-numerics} shows. The useful content is structural: the factor $2$
says localisation error and score error are equally damaging, and the appearance of $L_f$ says a
steep field tolerates less positional error, both of which are actionable when budgeting a
run. Note also that (a) is an assumption: it is delivered by using enough restarts, and
\cref{prop:no-global-guarantee} is exactly the statement that a single start cannot deliver it.
\end{remark}

\subsection{Dither: stated as a hypothesis, not a theorem}\label{sec:gap3-dither}
\NotationContext{field}

The second standard remedy is to add a decaying dither to the ring radius or to the centroid
command, in the spirit of extremum seeking \cite{krsticwang2000,khong2014}. Averaging theory
gives an asymptotic description of the resulting flow for sufficiently fast dither, but it does
not by itself produce a global convergence certificate for a fixed dither schedule on a fixed
multimodal field: the standard results are local around an averaged equilibrium, and matching an
escape-time argument to a decaying amplitude schedule requires assumptions on barrier widths and
depths that are not available here. Accordingly, dither is retained as a design hypothesis to be
assessed empirically, and no proposition is claimed for it in this document. What
\cref{thm:multistart} provides is independent of that choice: whatever escape mechanism is used,
the selection step is what converts a set of local outcomes into a certified global one.

\subsection{Numerical verification for Gap 3}\label{sec:gap3-numerics}

The test field is a three-component Gaussian mixture on $\Omega=[-6,6]^2$ with amplitudes
$(1.00,0.78,0.60)$, centres $(0,0)$, $(3.2,1.6)$, $(-2.8,1.9)$, and widths $(1.10,0.95,0.85)$.
Its peak values are $(1.0009,0.7852,0.6092)$, giving the value gap
$\Delta=0.2157$ used below.

\paragraph{Check 6: the bias exponent.}
$\norm{g_R(p)-\nabla f(p)}$ was measured at $p=(0.6,-0.4)$ over a decade of $R$ and a power law
was fitted. \cref{lem:estimator-bias} predicts slope $2$ for $n\geq4$;
\cref{lem:n3-bias} predicts slope $1$ for $n=3$.

\begin{center}
\begin{longtable}{@{}p{0.20\linewidth}p{0.34\linewidth}p{0.36\linewidth}@{}}
\toprule
$n$ & Predicted slope & Fitted slope \\
\midrule
\endhead
$3$ & $1$ (\cref{lem:n3-bias}) & $1.095$ \\
$4$ & $2$ (\cref{lem:estimator-bias}) & $1.989$ \\
$6$ & $2$ & $1.989$ \\
$8$ & $2$ & $1.990$ \\
\bottomrule
\end{longtable}
\end{center}

\paragraph{Check 7: the exact $n=3$ constant.}
The slope alone does not test \eqref{eq:n3-bias-vector}, only its order. At the same $p$ the
Hessian is $H_{11}=-0.4537$, $H_{12}=-0.1195$, $H_{22}=-0.5708$, so
\cref{lem:n3-bias} predicts
\begin{equation*}
    \frac{\norm{r_R}}{R}\to\frac14\sqrt{\left(H_{11}-H_{22}\right)^2+4H_{12}^2}
    =0.066520 .
\end{equation*}

\begin{center}
\begin{longtable}{@{}p{0.30\linewidth}p{0.30\linewidth}p{0.30\linewidth}@{}}
\toprule
Quantity at $R=10^{-4}$ & Predicted & Measured \\
\midrule
\endhead
Magnitude $\norm{r_R}/R$ & $0.066520$ & ratio $0.99998$ \\
Direction of $r_R$ & $\left(H_{11}-H_{22},\,-2H_{12}\right)$ & $0.015^\circ$ error \\
\bottomrule
\end{longtable}
\end{center}

Both the magnitude and the direction match, and the residual grows as $O(R)$ when $R$ is
increased, consistent with the $L_HR^2/3$ remainder of Step 6 of \cref{lem:n3-bias}. Because the
direction is a two-component prediction determined entirely by
\eqref{eq:n3-third-moment}, this check validates the moment computation of Steps 3--5 and not just
the scaling that Check 6 already established.

\paragraph{Check 8: how large the trapping set is.}
The ideal ascent flow \eqref{eq:ideal-ascent} was integrated from $48\,400$ initial conditions on
a uniform grid over $\Omega$ and each trajectory was attributed to the peak it reached.

\begin{center}
\begin{longtable}{@{}p{0.30\linewidth}p{0.24\linewidth}p{0.36\linewidth}@{}}
\toprule
Peak & Value & Fraction of initial conditions \\
\midrule
\endhead
$\ell_1$ (global) & $1.0009$ & $39.97\%$ \\
$\ell_2$ & $0.7852$ & $31.00\%$ \\
$\ell_3$ & $0.6092$ & $29.02\%$ \\
\bottomrule
\end{longtable}
\end{center}

A single start therefore fails on $60.03\%$ of initial conditions on this field. This is the
quantitative content of Step 7 of \cref{prop:no-global-guarantee}: the trapping set is not a thin
exceptional set but the majority of the domain, so restarts are a requirement rather than an
optimisation.

\paragraph{Check 9: the selection condition.}
\NotationContext{multistart}
Terminal outcomes were simulated for $M=3$ teams by perturbing true peak locations by
$\varepsilon_{\mathrm{loc}}$ and scores by $\delta_F$, drawn adversarially in sign, over a grid of
$(\varepsilon_{\mathrm{loc}},\delta_F)$ with $4\,000$ trials per grid point.

\begin{center}
\begin{longtable}{@{}p{0.44\linewidth}p{0.24\linewidth}p{0.22\linewidth}@{}}
\toprule
Region & Predicted & Measured accuracy \\
\midrule
\endhead
Inside \eqref{eq:multistart-condition} & always correct & $100\%$ \\
Widest uncertainty tested & no guarantee & $97.7\%$ \\
\bottomrule
\end{longtable}
\end{center}

Selection is exact everywhere \eqref{eq:multistart-condition} holds, as \cref{thm:multistart}
requires, and degrades only gradually outside it. This is the expected signature of a sufficient
worst-case condition: it is never violated where it applies, and it is conservative where it does
not, which is the point of \cref{rem:multistart-reading}.

\subsection{What the verification does not establish}\label{sec:gap3-limits}

Three limitations should be stated plainly, since each proof above is only as strong as its
hypotheses.

\begin{enumerate}[label=\arabic*.,nosep]
    \item The checks confirm the \emph{predicted scalings and directions}. They do not confirm
    that the constants $L_HR^2/3$, $\delta/(\kappa R)$ or $v_{\max}/(2\beta\kappa)$ are tight;
    they are upper bounds derived through triangle and Cauchy--Schwarz inequalities and are
    expected to be loose.
    \item \cref{prop:multisource-reduction} remains conditional on
    \eqref{eq:basin-certificate}, which the controller does not enforce
    (\cref{rem:gap1-conditional}). No numerical experiment can convert a conditional theorem into
    an unconditional one.
    \item Bounds of the form $\limsup_{t\to\infty}$, as in \eqref{eq:moving-ultimate-bound}, are
    asymptotic. A finite-horizon run showing a small tail error is supportive but strictly weaker
    evidence than the inequality claims, and a tail \emph{mean} below the bound is weaker still
    than a tail maximum.
\end{enumerate}

\begin{thebibliography}{9}

\bibitem{du2024}
B.~Du, K.~Qian, C.~Claudel, and D.~Sun,
``Multi-robot cooperative source seeking of a scalar field with a sign gradient-free algorithm,''
\emph{IEEE Transactions on Robotics}, 2024.

\bibitem{cortes2004}
J.~Cort\'es, S.~Mart\'inez, T.~Karata\c{s}, and F.~Bullo,
``Coverage control for mobile sensing networks,''
\emph{IEEE Transactions on Robotics and Automation}, vol.~20, no.~2, pp.~243--255, 2004.
doi: \href{https://doi.org/10.1109/TRA.2004.824698}{10.1109/TRA.2004.824698}.

\bibitem{michael2023}
E.~Michael, C.~Manzie, T.~A.~Wood, D.~Zelazo, and I.~Shames,
``Gradient free cooperative seeking of a moving source,''
\emph{Automatica}, vol.~152, art.~110948, 2023.
doi: \href{https://doi.org/10.1016/j.automatica.2023.110948}{10.1016/j.automatica.2023.110948}.

\bibitem{khong2014}
S.~Z.~Khong, Y.~Tan, C.~Manzie, and D.~Ne\v{s}i{\'c},
``Multi-agent source seeking via discrete-time extremum seeking control,''
\emph{Automatica}, vol.~50, no.~9, pp.~2312--2320, 2014.
doi: \href{https://doi.org/10.1016/j.automatica.2014.06.009}{10.1016/j.automatica.2014.06.009}.

\bibitem{krsticwang2000}
M.~Krsti{\'c} and H.-H.~Wang,
``Stability of extremum seeking feedback for general nonlinear dynamic systems,''
\emph{Automatica}, vol.~36, no.~4, pp.~595--601, 2000.
doi: \href{https://doi.org/10.1016/S0005-1098(99)00183-1}{10.1016/S0005-1098(99)00183-1}.

\end{thebibliography}

\end{document}
